\subsection{A Convex Family of Fairness Regularizers}
% \hh{Made this a separate section to emphasize the conceptual
% contribution (i.e. defining a family of fairness notions for
% regression)...}
Our formal definitions of fairness all measure how similarly a model
treats two similarly labeled instances, one from group $1$ and one
from group $2$.
% \hh{Shall we say one could consider taking into account not only the
% difference between y values, but also how close the explanatory
% variables are across the two instances?}
In particular, all of our definitions have a term for each
``cross-group'' \emph{pair} of instances/labels, weighted as a
function of $|\yi- \yj|$ and also by $|\vw\cdot \xi - \vw \cdot \xj|$.
For shorthand, we will refer to pairs of instances (one from each
group) as \emph{cross pairs}, and cross pairs with similar labels as
similar cross pairs.
%\jm{We need some shorthand for these things, I'm not sure if
%  this is what we should settle on but something seems useful.}  
Each of the below fairness definitions differs in precisely which
cross pair disparities can counteract one another. In one extreme
(individual fairness in Equation~\ref{eqn:innersq}), every cross pair
disparity increases the fairness penalty of a model. In the other
(group fairness in Equation~\ref{eqn:outersq}), making higher
predictions for the group $1$ instance of a similar cross pair can be
somewhat counterbalanced by making a higher prediction for the group
$2$ instance of a different similar cross pair. Our notions of
fairness for regression align closely to individual and group fairness
definitions for classification, both common threads in the fairness
literature.


\begin{remark}
  We assume the sensitive feature $\X_{d+1}$ is available to the
  learning procedure in one of two ways.  In the first setting, which
  we call the ``single model'' setting, we assume the algorithm builds a
  single linear model $\vw$ for all of $\X$ (over all but the
  sensitive features), but can measure the empirical fairness loss of
  $\vw$ using the sensitive feature. In the second setting, which we
  call the ``separate models'' setting, we allow the algorithm to build two
  distinct linear models $\vw_1, \vw_2$ for the two groups, $\vw_g$
  based on $S_g$, thus directly observing the sensitive feature when
  building these models.
\end{remark}


%\begin{remark}
We specialize the following fairness penalties for the single model
setting, but can easily extend them to the separate models setting, by
replacing $w$ with $w_g$ when applied to a member of group $g$.
%\end{remark}



\paragraph{Individual Fairness} The first fairness
penalty\footnote{For simplicity we define our fairness penalties on
  samples rather than on the underlying distribution $\cP$.  The
  analogous definitions with respect to the distribution can be
  derived by replacing sums with expectations.}  we propose is the
following:
\begin{align}
f_1(\vw, S) =\frac{1}{\n{1}\n{2}} \sum_{\substack{(\xi, \yi)\in S_1\\ (\xj,\yj)\in S_2}}
 %e^{-(\yi-\yj)^2} 
 d(\yi,\yj)\big(\vw\cdot \xi-\vw\cdot \xj\big)^2,\label{eqn:innersq}
\end{align}
for some fixed non-negative function $d$, which we assume is
decreasing in $|\yi - \yj|$ (see Section~\ref{sec:exp} for more
details).  Since $d(\yi, \yj)$ does not depend upon the decision
variables ($\vw$), one can treat these values as constants in an
optimization procedure for selecting $\vw$.

The penalty $f_1$ corresponds to \emph{\ind}; for every cross pair
$(\vx,y)\in S_1, (\vx', y')\in S_2$, a model $\vw$ is penalized for
how differently it treats $\vx$ and $\vx'$ (weighted by a function of
$|y - y'|$). No cancellation occurs: the penalty for overestimating
several of one group's labels cannot be mitigated by overestimating
several of the other group's labels. 


\paragraph{Group Fairness}
The second fairness penalty we propose is the following:
%
\begin{align}
f_2(\vw,S) =
\left( \frac{1}{\n{1}\n{2}} \sum_{\substack{(\xi,\yi)\in S_1\\ (\xj,\yj)\in S_2}} 
 d(\yi,\yj) \big(\vw\cdot \xi-\vw\cdot \xj\big)\right)^2.\label{eqn:outersq}
\end{align}

The penalty $f_2$ corresponds to \emph{\grp}: on average, the two
groups' instances should have similar labels (weighted by the nearness
of the labels of the instances). Unlike $f_1$, the penalty $f_2$
allows for \emph{compensation}:
%
%given
%$(\vx_1, y_1), (\vx_2, y_2)\in \X_1, (\vx'_1, y_1), (\vx'_2, y_2)\in
%\X_2$
%if the model predicts a higher value on $\vx_1$ than $\vx'_1$, and a
%higher value on $\vx'_2$ than $\vx_2$, then those terms counterbalance
%one another.
informally, if the model over-values some instances of group $1$
relative to group $2$ in similar cross pairs, it can compensate on
other similar cross-pairs by over-valuing those instances from group
$2$ relative to group $1$.

In both of the above formulations, for any cross pair
$(\xi,\yi) \in S_1$ and $(\xj, \yj) \in S_2$, any regressor $\vw$ will
have penalty that increases as
$\left|\vw \cdot \xi-\vw \cdot \xj\right|$ increases, weighted by
$d(\yi, \yj)$.  If the cross pair is similar ($\yi$ is close to $\yj$
and $d(\yi, \yj)$ is large), a regressor which makes very different
predictions for $\xi$ and $\xj$ will incur large loss. If the cross
pair is less similar ($\yi$ is far from $\yj$ and $d(\yi,\yj)$ is
smaller), there is less penalty for having a regressor for which
$\left|\vw \cdot \xi-\vw \cdot \xj\right|$ is large.

\paragraph{Hybrid notions of fairness}
Note that group and individual fairness correspond to two extremes: in
one extreme the fairness penalty considers each cross pair separately
and in the other one the fairness penalty considers all the cross
pairs together.  Mathematically one could define different notions of
fairness by grouping the cross pairs in different manners or even
restrict the fairness penalty only on a subset of cross pairs (called
\emph{bucketing}).  In particular, for binary labeled data (where
$\Y = \{-1,1\}$) one natural choice is to group the cross pairs based
on their labels.  This would result in the following definition of
fairness which we call~\hybrid:
\begin{align}
  f_3(\vw,S) &=
  \left(\sum_{\substack{(\xi,\yi) \in S_{1}\\ (\xj,\yj)\in S_{2}\\ \yi=\yj=1}}
  \frac{d(\yi,\yj) \big(\vw\cdot \xi-\vw\cdot \xj\big)}{n_{1,1} n_{2,1}}\right)^2 + 
  \left( \sum_{\substack{(\xi,\yi) \in S_{1}\\ (\xj,\yj)\in S_{2} \\ \yi=\yj=-1}}
  \frac{d(\yi,\yj) \big(\vw\cdot \xi-\vw\cdot \xj\big)}{n_{1,-1} n_{2,-1}}\right)^2,\label{eqn:betweensq}
\end{align}
where $n_{g,t}$ denotes the size of group $g\in\{1, 2\}$ with label
$t\in\{-1,1\}$ in the sample.  Intuitively, \hybrid~requires both
positive and both negatively labeled cross pairs to be treated
similarly in average over the two groups. Some compensation might
occur, but only amongst instances with the same label: over-valuing
positive instances from group $1$ can mitigate over-valuing positive
instances from group $2$, but does not mitigate under-valuing negative
instances from group $1$. These two terms can be weighted differently
for applications where the treatment of positive (or negative)
instances are more important.  See
Sections~\ref{sec:related}~and~\ref{sec:exp} for more details.

%As we show in our experiments in Section~\ref{sec:proxy}, \hybrid~can be used as a proxy for 
%enforcing similar rates of false positives and false negatives among the two groups.
%\sj{what is the relationship of this definition with~\citet{ZafarVGG17}?}

\subsection{Discussion of Our Notions of Fairness}

We now discuss several salient features of our fairness notions.

\paragraph{Why are these fairness notion different from accuracy?}
All of our fairness penalties are small for any perfect regressor (any
$\vw$ such that $\vw \cdot \vx = y$ for all $(\vx,y)\sim\cP$): for a
similar cross pair, $\yi \approx \yj$ and also
$\vw \cdot \xi \approx \vw \cdot \xj$ for a perfect regressor $\vw$.
Our fairness regularizers might then be interpreted as an unusual
proxy for standard accuracy rather than as fairness notions.  However,
perfect (linear or otherwise) regressors almost never exist in
practice; and between two models with similar accuracy, these
definitions bias a learning procedure towards those which have similar
treatment of similarly labeled instances from different groups.

\paragraph{What minimizes these penalties?}
We note that any constant regressor (any $\vw$ such that
$\vw \cdot \vx = c $ for all $\vx\in\X$ and some $c\in \R$) exactly
minimizes all of our fairness regularizers.  As we seen empirically,
this implies that as the fairness regularization factor $\lambda$
increases, we transition from an unfair model with minimum accuracy
loss to a constant (and therefore perfectly fair but trivial) model,
whose accuracy is the best any constant model can achieve.



%\sj{why these are good definitions for fairness}
%
%\hh{To be discussed: What does each notion of fairness mean/do
%  compared to accuracy alone? How did we come up with the formula? How
%  do they compare with individual and Group fairness notions in
%  previous work? ...}
%
%\hh{To be added: Shahin's argument as to why the last notion is
%  interesting in case of classification? Show that it reduces to
%  FP-FN.}
%  
%  \mj{Should insert some justification of why constant and perfect regressors
%  should both incur 0 fairness loss, or move this point to the discussion, 
%  perhaps in the definition discussion section.}

%\sj{is this true?perfect regressor might not have zero fairness loss for regression
%since we still have non-zero penalty of the form $e^{-(y-y')^2}$ for $y\ne y'$.}

%\begin{remark}
%  Consider a model $c : \bot \to \Y$ which produces outputs randomly
%  and independent of its input. Then, regularizer~\ref{eqn:innersq} is
%  non-zero for $c$ with high probability.
%  Regularizer~\ref{eqn:outersq} is also nonzero with high probability,
%  but is more tightly concentrated around $0$ for large $n$. \jm{I
%    don't think this remark is relevant anymore. We're not
%    axiomatizing fairness notions... though one could say that a
%    random classifier (at best?) can get the base rate of labelings}
%\end{remark}

% HH: commented out the line below
%We note that our fairness notions can be extended to \emph{linear}
%classifiers by replacing the loss $\ell_\cP(\vw)$ from mean-squared
%error of regression to a proper loss for classification.  \jm{do we do
%  this? I think we use the real-valued prediction... not sure if this
%  is distracting given we don't use binary outcomes in this work.}

